\section{Design de Experimentos para Modelagem e Avaliacao}

Esta etapa do projeto contempla o planejamento das fases de Modelagem e Avaliacao do CRISP-DM. O objetivo nao e executar os modelos, mas definir previamente o desenho experimental que sera usado na etapa seguinte. Dessa forma, a selecao futura da configuracao final sera baseada em um procedimento reprodutivel, com comparacao sistematica entre subconjuntos de atributos, arquiteturas de modelos, hiperparametros e metricas.

O fluxo experimental segue a hierarquia solicitada:

\begin{center}
\textbf{Features $\rightarrow$ Modelo $\rightarrow$ Hiperparametros $\rightarrow$ Metricas de Avaliacao}
\end{center}

A variavel alvo do problema e \texttt{serious}, tratada como uma tarefa de classificacao binaria. A base ja se encontra separada em treino e teste desde a etapa anterior, evitando vazamento de informacao entre os conjuntos. No conjunto de treino, a distribuicao da classe alvo e moderadamente desbalanceada, com aproximadamente 62,3\% dos registros na classe \texttt{1} e 37,7\% na classe \texttt{2}. Por isso, a metrica primaria de selecao planejada e o \textit{F1-score macro}, complementado por acuracia balanceada e demais metricas de diagnostico.

\subsection{Subconjuntos de Atributos}

Foram definidos tres subconjuntos incrementais de atributos. Essa organizacao permite avaliar se a inclusao de informacoes estruturadas de medicamentos e de atributos textuais agregados melhora o desempenho dos modelos.

\begin{table}[htbp]
\centering
\caption{Subconjuntos de atributos planejados para o Design de Experimentos.}
\label{tab:feature_sets_doe}
\begin{tabular}{p{0.22\textwidth} p{0.68\textwidth}}
\hline
\textbf{Feature set} & \textbf{Descricao} \\
\hline
\texttt{F1\_basico\_clinico} & Atributos demograficos e clinicos tratados, incluindo pais de ocorrencia, sexo do paciente, faixa etaria calculada, indicador de imputacao e indicador de ocorrencia nos Estados Unidos. \\
\texttt{F2\_medicamentos\_estruturado} & Conjunto \texttt{F1} acrescido de contagens e indicadores estruturados de medicamentos, como quantidade de medicamentos, numero de substancias ativas, tipos de medicamento e indicadores binarios derivados. \\
\texttt{F3\_completo\_com\_texto} & Conjunto \texttt{F2} acrescido dos campos textuais agregados de substancia ativa e produto medicinal, vetorizados no pipeline de modelagem. \\
\hline
\end{tabular}
\end{table}

\subsection{Modelos e Hiperparametros}

O planejamento considera modelos com diferentes niveis de complexidade e interpretabilidade. A Regressao Logistica funciona como baseline linear; a Arvore de Decisao como baseline nao linear simples; Random Forest e Gradient Boosting representam ensembles; e SVM permite avaliar classificadores baseados em margem. Para cada modelo sera usada uma busca sistematica de hiperparametros por \texttt{GridSearchCV} ou \texttt{RandomizedSearchCV}, conforme o tamanho do espaco de busca.

\begin{table}[htbp]
\centering
\caption{Modelos e hiperparametros planejados.}
\label{tab:modelos_hiperparametros_doe}
\begin{tabular}{p{0.22\textwidth} p{0.18\textwidth} p{0.48\textwidth}}
\hline
\textbf{Modelo} & \textbf{Busca} & \textbf{Hiperparametros planejados} \\
\hline
Regressao Logistica & \texttt{GridSearchCV} & \texttt{C}, \texttt{penalty}, \texttt{solver}. \\
Arvore de Decisao & \texttt{GridSearchCV} & \texttt{max\_depth}, \texttt{min\_samples\_leaf}, \texttt{criterion}. \\
Random Forest & \texttt{RandomizedSearchCV} & \texttt{n\_estimators}, \texttt{max\_depth}, \texttt{min\_samples\_leaf}, \texttt{max\_features}. \\
Gradient Boosting & \texttt{RandomizedSearchCV} & \texttt{n\_estimators}, \texttt{learning\_rate}, \texttt{max\_depth}, \texttt{subsample}. \\
SVM & \texttt{GridSearchCV} & \texttt{C}, \texttt{kernel}, \texttt{gamma}. \\
\hline
\end{tabular}
\end{table}

\subsection{Validacao Cruzada e Metricas}

A validacao cruzada planejada sera estratificada, usando \texttt{StratifiedKFold} com cinco particoes, embaralhamento e semente fixa. A estratificacao preserva a proporcao da variavel alvo em cada particao da validacao cruzada, reduzindo a variabilidade provocada pelo desbalanceamento moderado das classes.

As metricas serao documentadas para treinamento, validacao cruzada e teste. O uso de metricas nos tres niveis permite identificar sobreajuste, estabilidade da configuracao e capacidade de generalizacao.

\begin{table}[htbp]
\centering
\caption{Metricas planejadas para avaliacao dos experimentos.}
\label{tab:metricas_doe}
\begin{tabular}{p{0.30\textwidth} p{0.58\textwidth}}
\hline
\textbf{Metrica} & \textbf{Finalidade} \\
\hline
\texttt{accuracy} & Medir a proporcao total de classificacoes corretas. \\
\texttt{balanced\_accuracy} & Avaliar desempenho medio por classe, reduzindo o efeito do desbalanceamento. \\
\texttt{precision\_macro} & Medir a precisao media entre as classes. \\
\texttt{recall\_macro} & Medir a revocacao media entre as classes. \\
\texttt{f1\_macro} & Metrica primaria de selecao, equilibrando precisao e revocacao por classe. \\
\texttt{roc\_auc} & Avaliar capacidade de discriminacao quando o modelo disponibilizar probabilidades ou funcao de decisao. \\
\hline
\end{tabular}
\end{table}

\subsection{Estrutura de Resultados Esperada}

Para cada combinacao de subconjunto de atributos, modelo e hiperparametros, a etapa futura de execucao devera registrar: melhores hiperparametros, metricas de treino, media e desvio-padrao das metricas na validacao cruzada, e metricas finais no conjunto de teste. A configuracao final sera escolhida pela maior media de \texttt{f1\_macro} na validacao cruzada. Em caso de empate pratico, serao considerados, nesta ordem, menor variancia do \texttt{f1\_macro}, maior \texttt{balanced\_accuracy}, menor complexidade do modelo e desempenho consistente no teste.

Assim, esta entrega deixa preparado o Design de Experimentos, mas nao realiza o ajuste dos modelos. A execucao dos objetos \texttt{GridSearchCV} e \texttt{RandomizedSearchCV}, bem como o preenchimento das metricas finais, deve ocorrer apenas na etapa seguinte.
